\section{Выводы}

В ходе работы я пришел к следующим выводам:
\begin{itemize}
    \item \textbf{Эффективность профилирования:} Использование утилит \texttt{gprof} и \texttt{Valgrind} позволило наглядно увидеть распределение ресурсов внутри программы. Без профилирования выявление скрытых задержек, связанных с избыточным копированием объектов, было бы затруднительным.
    
    \item \textbf{Оптимизация перемещения (Move semantics):} Внедрение \texttt{std::move} в критических узлах программы (вставка элементов) показало существенный результат: количество аллокаций памяти снизилось на 9\%, что напрямую уменьшает нагрузку на системный менеджер памяти и снижает риск фрагментации.
    
    \item \textbf{Безопасность и стек:} Переход от рекурсивной очистки дерева к итеративной реализации (с использованием очереди) позволил гарантировать стабильность программы. Это критически важно при работе с потенциально глубокими префиксными деревьями, где стандартный стек вызовов может стать «узким местом».
    
    \item \textbf{Анализ «хотспотов»:} Сравнение профилей версий v1 и v2 показало, что после устранения системного шума (копирования строк) программа стала тратить время преимущественно на целевую логику — побитовый поиск в дереве.
    
    \item \textbf{Практическая ценность:} Исследование подтвердило, что даже алгоритмически оптимальная структура (PATRICIA) может работать неэффективно при неаккуратном управлении ресурсами языка C++. Навык низкоуровневой оптимизации на основе данных профилировщика является ключевым для разработки высокопроизводительных систем.
\end{itemize}

Данный опыт закрепил понимание жизненного цикла объектов в памяти и научил принимать инженерные решения не только на основе теоретической сложности алгоритма, но и на основе реальных показателей потребления ресурсов.

\pagebreak